\section{Relevance to Proof Systems and Limitations}
\label{sec:proof-systems}

The experiments in this paper isolate arbitrary single-point evaluation of a
univariate polynomial already stored in a chosen representation.  This
section explains where such a primitive appears in existing proof-system
architectures, which parts of the optimization can transfer directly, and
which conclusions would require end-to-end integration experiments.

\subsection{Where arbitrary-point evaluations appear}

Transparent proof systems based on Reed--Solomon proximity testing commonly
move between polynomial coefficients and evaluations on structured domains.
They also evaluate committed polynomials at verifier-derived challenge
points that need not lie on the original commitment domain.

RISC Zero provides a concrete public example of this pattern.  Its STARK
documentation describes conversion of trace columns to polynomial
coefficients with an inverse NTT, Reed--Solomon evaluation on an expanded
domain, DEEP-ALI checks at an out-of-domain point \(z\), and a subsequent
FRI low-degree argument \cite{risc0starkbyhand,risc0sequence}.

In the DEEP-ALI stage, the prover supplies evaluations of trace polynomials
at points of the form
\[
P_i(\omega^n z),
\]
where \(z\) is a verifier challenge and the shifts depend on the constraint
system \cite{risc0sequence}.  This is the class of operation for which an
efficient arbitrary-point evaluator can potentially matter.

The relevance is nevertheless conditional on representation.  Existing
STARK implementations are typically organized around evaluation-domain
vectors because FFT/NTT-based encoding, pointwise constraint operations,
Merkle commitments, and FRI naturally consume such data.  Replacing only
one evaluation routine without considering the representation of the
surrounding pipeline may therefore introduce conversion costs that dominate
the local gain.

\subsection{Relation to FRI}

FRI proves that a committed vector is close to evaluations of a low-degree
polynomial.  Its commit phase recursively folds evaluations into smaller
commitments, while the query phase opens a small number of positions from
the committed layers \cite{risc0fri}.

The Boolean-kernel fold studied in this paper,
\[
F_t(a,b)=t(a+b)+b,
\]
should not be identified with a standard FRI fold merely because both are
binary reductions.  They act on different representations and satisfy
different protocol roles.

The present result instead motivates the following systems question:

\begin{quote}
Can a prover retain a Boolean-kernel representation across enough of the
polynomial pipeline that the fast arbitrary-point evaluator is useful
without paying an offsetting conversion cost before FRI or commitment?
\end{quote}

Answering this requires protocol-level design and benchmarking.  The current
paper establishes the performance of the local representation primitive;
it does not establish a new FRI protocol.

\subsection{Transferable low-level optimizations}

Some optimizations presented in \cref{sec:implementation} are not specific
to the Boolean-kernel basis.

The specialized Goldilocks reduction exploits only
\[
p=2^{64}-2^{32}+1.
\]
It can therefore benefit other Goldilocks workloads when the same arithmetic
pattern occurs.

This is relevant because Goldilocks has already been used specifically to
obtain efficient CPU arithmetic in FRI-based proof systems.  Plonky2, for
example, was designed around FRI and the Goldilocks field precisely because
the modulus permits efficient \(64\)-bit arithmetic on modern processors
\cite{plonky2intro}.

Likewise, the experiments on allocation reuse, worker-local scratch space,
bounds-check elimination, and task granularity are general implementation
lessons for prover kernels.  Their exact gains are workload-dependent, but
they do not rely on a change to the cryptographic protocol.

By contrast, the speedup arising from the Boolean-kernel representation
itself transfers only if the relevant polynomial is actually available in
that representation.

\subsection{Effect on an existing prover}

Suppose an existing prover currently represents a polynomial by monomial
coefficients or by evaluations on a roots-of-unity domain.

There are three distinct integration cases.

\paragraph{Case 1: local arithmetic replacement only.}
If the prover retains its existing representation and proof format but
adopts only the Goldilocks arithmetic optimizations, then the cryptographic
protocol is unchanged.  Any speedup is limited to the affected arithmetic
kernels.

\paragraph{Case 2: internal representation change with identical proof
interface.}
If the prover internally stores some polynomials in Boolean-kernel form but
converts back before commitment or proof serialization, then an external
verifier need not change.  The benefit is positive only if
\[
T_{\mathrm{saved\ evaluations}}
>
T_{\mathrm{basis\ conversions}}
+
T_{\mathrm{integration\ overhead}}.
\]

\paragraph{Case 3: protocol-visible representation change.}
If commitments, openings, FRI layers, transcript messages, or verifier
equations themselves change, then the proof format changes.  Existing
verifiers cannot automatically accept the new proof format.

This distinction is important when discussing systems that already have
deployed verification infrastructure.

\subsection{On-chain verification}

A prover optimization does not by itself reduce on-chain verification gas.

For example, RISC Zero publishes Solidity verifier interfaces for receipts
and Ethereum integration contracts \cite{risc0ethverifier}.  These contracts
consume a particular proof or receipt format.

If a faster prover produces exactly the same proof bytes and verification
relation, then:
\[
\boxed{
\text{prover time may decrease while on-chain verification remains unchanged}.
}
\]

In this case no smart-contract change is required, but gas usage is not
expected to improve merely because the prover computed the proof faster.

If instead the Boolean-kernel representation changes the proof object or
verification equations, then deployed contracts would need a compatible new
verifier implementation, a routing/versioning mechanism, or a migration
path.  Existing RISC Zero Ethereum infrastructure, for example, uses
verifier selectors and routing between verifier implementations, illustrating
that verifier parameters and proof formats are versioned objects
\cite{risc0router}.

Accordingly, the present benchmark should be viewed primarily as a
\emph{prover-side} result.

\subsection{When the benchmark is likely to matter}

The representation is most promising in workloads satisfying several of
the following conditions:

\begin{enumerate}
    \item large polynomials are evaluated repeatedly at verifier-derived
    off-domain points;
    \item the polynomial can remain in Boolean-kernel representation for
    more than one operation;
    \item basis conversion can be amortized or eliminated;
    \item the prover is CPU-bound in field arithmetic or memory movement
    rather than dominated by hashing or external I/O;
    \item the target field admits efficient realization of the fold
    arithmetic.
\end{enumerate}

Conversely, the result is less likely to affect workloads in which the
polynomial is used almost exclusively through evaluations on one fixed FFT
domain or where a single arbitrary-point evaluation is surrounded by much
larger hashing, commitment, or encoding costs.

\subsection{Representation conversion}

The monomial-to-kernel conversion described in
\cref{sec:representations} is a Boolean Möbius transform and costs
\[
O(N\log N)
\]
field operations.

This cost is not included in the single-point evaluation timing because
the benchmark assumes the input is already represented in the basis being
measured.

For a one-off query, conversion can easily outweigh the evaluation gain.
The basis is therefore most interesting when either:

\begin{enumerate}
    \item the witness or polynomial is generated directly in kernel
    coefficients;
    \item several subsequent operations consume the kernel representation;
    \item conversion is fused with another prover stage;
    \item or the representation replaces another transform already present
    in the prover.
\end{enumerate}

An end-to-end evaluation must account for the full representation lifetime,
not only the isolated query.

\subsection{Limitations of the present experiments}

The current study has several explicit limitations.

\paragraph{Single processor family.}
The headline measurements are from an Apple M1 system.  Cache sizes,
instruction scheduling, multiplication latency, memory bandwidth, and the
cost of Rayon scheduling differ on x86-64 servers and other ARM systems.

\paragraph{Single base field.}
The final measurements use Goldilocks.  The Boolean-kernel algebra is
field-independent, but the constant factors of the implementation are not.
A field such as Mersenne-31 may produce a different balance between
arithmetic and memory costs.

\paragraph{CPU implementation.}
No GPU, FPGA, or wide-SIMD implementation is evaluated in this paper.

\paragraph{Single-point evaluation focus.}
The paper does not yet benchmark all operations needed by a proof system.
In particular, full-domain encoding, low-degree extension, commitment,
basis conversion, and FRI integration are separate workloads.

\paragraph{No end-to-end prover measurement.}
The measured
\[
1.337\times
\]
same-run speed advantage over the parallel monomial baseline applies to the
single-point evaluator only.  It is not an end-to-end STARK or PCS speedup.

\paragraph{Run-to-run variability.}
The best kernel-only optimization run reached
\[
892.57\,\mu\mathrm{s},
\]
while the final same-run cross-basis sweep measured
\[
1.0560\,\mathrm{ms}.
\]
The paper therefore uses same-run measurements when constructing
cross-basis ratios and reports the lower number only as a best observed
implementation result.

\paragraph{Lagrange implementation scope.}
The roots-of-unity Lagrange benchmark uses the direct barycentric formula
with batch inversion for an arbitrary point.  It is not intended to replace
FFT/NTT evaluation when the requested points form a structured domain.

\subsection{Next system-level experiments}

The benchmark naturally suggests the following next experiments:

\begin{enumerate}
    \item measure monomial-to-kernel and kernel-to-monomial conversion
    alongside single-point evaluation;
    \item measure repeated queries to determine the query count at which
    conversion amortizes;
    \item integrate the evaluator into a DEEP-style off-domain evaluation
    stage of an existing STARK implementation;
    \item benchmark a kernel-native route to FRI commitment and folding;
    \item repeat the study on x86-64, Mersenne-31, and additional processor
    architectures;
    \item measure the fraction of total prover time attributable to the
    operation before making an end-to-end performance claim.
\end{enumerate}

The principal open engineering question is therefore no longer whether the
Boolean-kernel representation supports a fast arbitrary-point evaluation.
The experiments establish that it does.  The remaining question is whether
a complete proof-system pipeline can preserve the representation long enough
for this local advantage to survive at system scale.
